\documentclass[11pt,a4paper]{article}

\usepackage[T1]{fontenc}
\usepackage[utf8]{inputenc}
\usepackage[italian]{babel}
\usepackage[a4paper,margin=2.3cm]{geometry}
\usepackage{xcolor}
\usepackage{booktabs}
\usepackage{listings}
\usepackage[hidelinks]{hyperref}

\setlength{\emergencystretch}{3em}

\definecolor{codebackground}{RGB}{246,248,250}
\definecolor{codeframe}{RGB}{208,215,222}

\lstdefinestyle{shell}{
  language=bash,
  basicstyle=\ttfamily\small,
  backgroundcolor=\color{codebackground},
  frame=single,
  rulecolor=\color{codeframe},
  breaklines=true,
  columns=fullflexible,
  keepspaces=true,
  showstringspaces=false,
  upquote=true
}

\title{Cooling Solver\\Compilazione su Leonardo}
\author{}
\date{Aggiornata al \today}

\begin{document}

\maketitle
%\tableofcontents

\section{Scopo}

Questa guida descrive la compilazione manuale del Cooling Solver sul cluster
Leonardo. Sono mantenuti due profili distinti:

\begin{description}
  \item[Confronto tra implementazioni] produce file HDF5 da validare e usa
  opzioni floating-point conservative per ridurre le differenze numeriche;
  \item[Benchmark dello speed-up] elimina completamente HDF5 e abilita le
  ottimizzazioni più aggressive appropriate per la CPU Ice Lake e la GPU
  NVIDIA A100 di Leonardo.
\end{description}

I due profili producono eseguibili con nomi diversi, per evitare di usare per
errore una build di validazione durante un benchmark. Tutti i comandi devono
essere lanciati dalla directory principale del progetto.

\section{Ambiente software}

\subsection{Moduli comuni}

Iniziare da un ambiente riproducibile e caricare NVHPC:

\begin{lstlisting}[style=shell]
module purge
module load nvhpc

mkdir -p install/bin
\end{lstlisting}

\subsection{Modulo aggiuntivo per il confronto HDF5}

Solo per il profilo di confronto è necessario caricare HDF5:

\begin{lstlisting}[style=shell]
module load hdf5/1.14.3--gcc--12.2.0-spack0.22

HDF5_PREFIX="$(dirname "$(dirname "$(command -v h5c++)")")"
HDF5_INCLUDE="${HDF5_PREFIX}/include"
HDF5_LIB="${HDF5_PREFIX}/lib"

\end{lstlisting}


\section{Profilo 1: confronto tra implementazioni}

\subsection{Obiettivo del profilo}

Questo profilo serve a confrontare i campi prodotti dalle quattro
implementazioni. Tutti gli eseguibili:

\begin{itemize}
  \item sono compilati con lo stesso compilatore \texttt{nvc++};
  \item supportano la scrittura HDF5;
  \item usano doppia precisione e opzioni floating-point conservative;
  \item mantengono \texttt{-O3}, perché il confronto deve comunque essere
  rappresentativo di una normale build ottimizzata.
\end{itemize}

Queste opzioni riducono alcune fonti di differenza, ma non garantiscono
uguaglianza bit-a-bit: l'ordine delle operazioni e delle riduzioni cambia tra
seriale, OpenMP, OpenACC multicore e GPU.

\subsection{CPU seriale con HDF5}

\begin{lstlisting}[style=shell]
nvc++ -O3 -std=c++17 \
  -Kieee \
  -Mnofma \
  -DUSE_HDF5=1 \
  -I"${HDF5_INCLUDE}" \
  src/cpp/cooling.cpp \
  -o install/bin/cooling_cpp_compare \
  -L"${HDF5_LIB}" \
  -Wl,-rpath,"${HDF5_LIB}" \
  -lhdf5_cpp -lhdf5 -lz -ldl -lm -lstdc++fs
\end{lstlisting}

\subsection{CPU OpenMP con HDF5}

\begin{lstlisting}[style=shell]
nvc++ -O3 -std=c++17 \
  -mp \
  -Kieee \
  -Mnofma \
  -DUSE_HDF5=1 \
  -I"${HDF5_INCLUDE}" \
  src/cpp/cooling_openmp.cpp \
  -o install/bin/cooling_omp_compare \
  -L"${HDF5_LIB}" \
  -Wl,-rpath,"${HDF5_LIB}" \
  -lhdf5_cpp -lhdf5 -lz -ldl -lm -lstdc++fs
\end{lstlisting}

\subsection{CPU OpenACC multicore con HDF5}

\begin{lstlisting}[style=shell]
nvc++ -O3 -std=c++17 \
  -acc=multicore \
  -Minfo=accel \
  -Kieee \
  -Mnofma \
  -DUSE_HDF5=1 \
  -I"${HDF5_INCLUDE}" \
  src/cpp/cooling_openacc.cpp \
  -o install/bin/cooling_acc_cpu_compare \
  -L"${HDF5_LIB}" \
  -Wl,-rpath,"${HDF5_LIB}" \
  -lhdf5_cpp -lhdf5 -lz -ldl -lm -lstdc++fs
\end{lstlisting}

\subsection{GPU OpenACC A100 con HDF5}

\begin{lstlisting}[style=shell]
nvc++ -O3 -std=c++17 \
  -acc=gpu \
  -gpu=cc80,nofma \
  -Minfo=accel \
  -Kieee \
  -Mnofma \
  -DUSE_HDF5=1 \
  -I"${HDF5_INCLUDE}" \
  src/cpp/cooling_openacc.cpp \
  -o install/bin/cooling_acc_gpu_compare \
  -L"${HDF5_LIB}" \
  -Wl,-rpath,"${HDF5_LIB}" \
  -lhdf5_cpp -lhdf5 -lz -ldl -lm -lstdc++fs
\end{lstlisting}

In questo caso \texttt{nofma} viene specificato anche dentro \texttt{-gpu},
affinché la contrazione FMA sia disabilitata esplicitamente nel codice
eseguito sulla GPU.


\section{Profilo 2: benchmark dello speed-up}

\subsection{Obiettivo del profilo}

Queste build servono esclusivamente a misurare i tempi. Non definiscono
\texttt{USE\_HDF5}, non includono header HDF5 e non collegano librerie HDF5.
In esecuzione bisogna passare \texttt{none} come secondo argomento e
\texttt{0} come intervallo di output.

\subsection{CPU seriale ottimizzata}

\begin{lstlisting}[style=shell]
nvc++ -O3 -fast -tp=icelake -std=c++17 \
  src/cpp/cooling.cpp \
  -o install/bin/cooling_cpp_bench \
  -lstdc++fs
\end{lstlisting}

\subsection{CPU OpenMP ottimizzata}

\begin{lstlisting}[style=shell]
nvc++ -O3 -fast -tp=icelake -std=c++17 \
  -mp \
  src/cpp/cooling_openmp.cpp \
  -o install/bin/cooling_omp_bench \
  -lstdc++fs
\end{lstlisting}

\subsection{CPU OpenACC multicore ottimizzata}

\begin{lstlisting}[style=shell]
nvc++ -O3 -fast -tp=icelake -std=c++17 \
  -acc=multicore \
  -Minfo=accel \
  src/cpp/cooling_openacc.cpp \
  -o install/bin/cooling_acc_cpu_bench \
  -lstdc++fs
\end{lstlisting}

\subsection{GPU OpenACC A100 ottimizzata}

\begin{lstlisting}[style=shell]
nvc++ -O3 -fast -std=c++17 \
  -acc=gpu \
  -gpu=cc80,fastmath \
  -Minfo=accel \
  src/cpp/cooling_openacc.cpp \
  -o install/bin/cooling_acc_gpu_bench \
  -lstdc++fs
\end{lstlisting}


\section{Spiegazione dei flag}

\subsection{Flag comuni}

\begin{description}
  \item[\texttt{-O3}] abilita un livello elevato di ottimizzazione, incluse
  trasformazioni aggressive dei loop;
  \item[\texttt{-std=c++17}] seleziona lo standard C++17 richiesto dal codice;
  \item[\texttt{-lstdc++fs}] collega il supporto GNU per
  \texttt{std::filesystem}; è mantenuto per compatibilità con lo stack del
  cluster;
  \item[\texttt{-o file}] stabilisce il nome dell'eseguibile prodotto.
\end{description}

\subsection{Flag di parallelizzazione}

\begin{description}
  \item[\texttt{-mp}] abilita le direttive OpenMP e collega il runtime OpenMP
  di NVHPC;
  \item[\texttt{-acc=multicore}] traduce le direttive OpenACC per eseguirle sui
  core della CPU;
  \item[\texttt{-acc=gpu}] traduce le direttive OpenACC generando kernel per
  GPU NVIDIA;
  \item[\texttt{-gpu=cc80}] genera codice per compute capability 8.0, cioè
  l'architettura delle GPU NVIDIA A100 dei nodi Booster di Leonardo;
  \item[\texttt{-Minfo=accel}] stampa a compilazione conclusa quali regioni e
  loop sono stati trasformati per il target parallelo. Non introduce overhead
  durante l'esecuzione.
\end{description}

\subsection{Flag del profilo di confronto}

\begin{description}
  \item[\texttt{-Kieee}] richiede un comportamento floating-point più rigoroso
  rispetto alle regole IEEE e limita alcune trasformazioni algebriche;
  \item[\texttt{-Mnofma}] impedisce al compilatore host di contrarre una
  moltiplicazione e una somma in una singola FMA;
  \item[\texttt{-gpu=cc80,nofma}] seleziona A100 e disabilita FMA anche nel
  codice dispositivo;
  \item[\texttt{-DUSE\_HDF5=1}] attiva nel sorgente le sezioni protette da
  \texttt{\#ifdef USE\_HDF5};
  \item[\texttt{-I}] aggiunge la directory degli header HDF5;
  \item[\texttt{-L}] aggiunge la directory delle librerie HDF5 al linker;
  \item[\texttt{-Wl,-rpath,...}] registra nell'eseguibile il percorso delle
  librerie HDF5, così il loader può ritrovarle durante il job;
  \item[\texttt{-lhdf5\_cpp -lhdf5}] collega rispettivamente le API C++ e C di
  HDF5;
  \item[\texttt{-lz -ldl -lm}] collega le dipendenze usate dallo stack HDF5:
  compressione zlib, caricamento dinamico e libreria matematica.
\end{description}

\subsection{Flag del profilo benchmark}

\begin{description}
  \item[\texttt{-fast}] abilita il gruppo di ottimizzazioni NVHPC consigliate
  per le prestazioni. Il contenuto preciso del gruppo può cambiare tra versioni
  del compilatore;
  \item[\texttt{-tp=icelake}] genera codice host specifico per le CPU Intel Ice
  Lake dei nodi Booster, consentendo l'uso delle relative istruzioni vettoriali;
  \item[\texttt{-gpu=cc80,fastmath}] seleziona A100 e usa implementazioni
  matematiche GPU più rapide e meno conservative. Può modificare gli ultimi bit
  dei risultati e non deve essere usato per stabilire la tolleranza del
  validator.
\end{description}

Nel profilo benchmark sono intenzionalmente assenti \texttt{-Kieee},
\texttt{-Mnofma}, \texttt{-DUSE\_HDF5}, i percorsi HDF5 e tutte le librerie
HDF5. In questo modo la misura non comprende scrittura dei campi né vincoli
floating-point introdotti esclusivamente per il debug numerico.


\section{Selezione del profilo negli script SLURM}

Lo script \texttt{submit\_cpp.sh} può selezionare l'eseguibile seriale corretto
in base a un argomento passato a \texttt{sbatch}. Poiché i lanci di benchmark
saranno più frequenti, \texttt{benchmark} è la modalità predefinita: se non si
specifica alcun argomento viene eseguito \texttt{cooling\_cpp\_bench} senza
output HDF5.

Prima di inviare il primo job bisogna creare le directory usate per log e
risultati. In particolare, \texttt{logs/} deve esistere prima di
\texttt{sbatch}, perché Slurm apre i file di standard output e standard error
prima di eseguire lo script:

\begin{lstlisting}[style=shell]
mkdir -p logs output
\end{lstlisting}

I possibili comandi di lancio sono:

\begin{center}
\begin{tabular}{@{}lll@{}}
\toprule
Comando & Profilo & Eseguibile \\
\midrule
\texttt{sbatch submit\_cpp.sh} & benchmark & \texttt{cooling\_cpp\_bench} \\
\texttt{sbatch submit\_cpp.sh benchmark} & benchmark & \texttt{cooling\_cpp\_bench} \\
\texttt{sbatch submit\_cpp.sh compare} & confronto & \texttt{cooling\_cpp\_compare} \\
\bottomrule
\end{tabular}
\end{center}

Nel profilo \texttt{benchmark} lo script passa \texttt{none} al posto del file
HDF5 e usa un intervallo di output pari a zero. Nel profilo \texttt{compare}
scrive invece \texttt{output/Cooling\_cpp.h5} ogni 20 time step.

In modo analogo possono essere chiamati gli altri script per i vari casi: 
\texttt{submit\_omp.sh}, \texttt{submit\_acc\_cpu.sh}, e \texttt{submit\_acc\_gpu.sh}.

\section{Installazione dell'ambiente Python e validazione}

Il validator \texttt{tools/validate\_cooling\_h5.py} usa i moduli Python
\texttt{h5py} e \texttt{numpy}. Il modulo \texttt{python/3.11.7} di Leonardo
fornisce l'interprete, ma non garantisce che questi pacchetti siano installati.
Il progetto usa quindi un virtual environment chiamato
\texttt{cooling\_venv}, collocato nella directory principale del repository.
L'installazione deve essere eseguita una volta sola. 

Per creare l'environment si invia il job dalla directory principale del progetto:

\begin{lstlisting}[style=shell]
sbatch submit_install_pyenv.sh
\end{lstlisting}

Con le tolleranze definite nello script, lanciare:

\begin{lstlisting}[style=shell]
sbatch submit_validate.sh
\end{lstlisting}


Durante il confronto, \texttt{/step} viene verificato esattamente, mentre per
\texttt{/field} viene applicata la condizione

\[
  \left|u_{\mathrm{cand}}-u_{\mathrm{ref}}\right|
  \leq
  \mathrm{ATOL} + \mathrm{RTOL}\,
  \left|u_{\mathrm{ref}}\right|.
\]



\end{document}
